\documentclass{article}
\usepackage{graphicx} % Required for inserting images
\usepackage{xcolor} % Must be loaded before soul
\usepackage{soul} % For highlighting text
\usepackage[utf8]{inputenc}
\usepackage{newunicodechar}
\usepackage{tikz}
\usetikzlibrary{positioning}
\usepackage[utf8]{inputenc}
\usepackage[T1]{fontenc}
\usepackage{fontspec}           % For Greek text support with XeLaTeX/LuaLaTeX
\usepackage{polyglossia}        % Multilingual support
\setmainlanguage{english}
\setotherlanguage{greek}

\usepackage{geometry}
\geometry{margin=1in}

\usepackage{csquotes}           % Block quotes
\usepackage{enumitem}           % Better lists
\usetikzlibrary{arrows.meta, positioning, shapes.geometric}

\usepackage{fancyvrb}           % Verbatim environments for ASCII diagrams
\usepackage{setspace}           % Line spacing



% --- Greek support (XeLaTeX) ---
\usepackage{fontspec}
\newfontfamily\greekfont{Gentium Plus} % good polytonic Greek; try Times New Roman if needed
\newcommand{\gk}[1]{{\greekfont #1}}

\newunicodechar{₁}{$_1$}

\usepackage[style=mla]{biblatex} % MLA Citation Style for Bibliography
\addbibresource{references.bib} %Add References page for in text citation
\usepackage[colorlinks=true, linkcolor=blue, urlcolor=blue, citecolor=blue]{hyperref} %enable hyperlinks to be colored and interactable
\usepackage{changepage} %create individual blocks of text that can be altered apart from general document.

\newcommand{\ra}{\ensuremath{\rightarrow}}

\newcommand{\inlinenote}[1]{\par\noindent
  \fcolorbox{orange}{orange!20}{\parbox{0.97\linewidth}{#1}}\par\noindent}

\begin{document}

\section{$M_0 \rightarrow A_1$: The Actualization of Perception}

\subsection{The Preconditions of $A_1$: Sensible Object, Sense-Organ, and Medium}

$A_1$ is the completed perceptual act --- the joint actualization, licensed by $A_0$'s motion-time horizon, of rational souls, sensible objects actively exercising their sensible powers, sense-organs holding their first actuality as capacities, and a temporal-narrative manifold within which the actualization chain proceeds --- together with the residual trace the act deposits at the rational soul. The perceptual motion $M_0 \rightarrow A_1$ is the \textit{kinēsis} (\gk{κίνησις} `motion') whose terminus is $A_1$ and whose unmoved originator is the sensible object's active potency engaging at $A_0$'s motion-time substrate. The kinetic-roles structure inherited from $A_0$ holds throughout: the unmoved originator is the sensible object exercising its proper quality in actuality; the moved mover is the sense faculty receiving the sensible form (\textit{eidos}, \gk{εἶδος} `form') without the matter (\textit{hylē}, \gk{ὕλη} `matter'); the animal via the sense-organ is that which is moved.\footnote{The three-factor schema is given at \textit{De Anima} III.10, 433b13--18: ``these three are required: the mover, the means by which it moves, and the moved\ldots\ thus there are three things: the mover, the means, and the moved.'' Aristotle reiterates the schema's coupling to \textit{phantasia} at 433b27--30: ``inasmuch as an animal is capable of appetite it is capable of self-movement; it is not capable of appetite without possessing imagination; and all imagination is either calculative or sensitive.'' §1.1 established the schema's primary application to locomotion at the \textit{orektikon}-stage of the chain; the present section deploys its analogical extension to the perceptual stage, where the sense-organ functions as the moved mover and the sensible object as the unmoved originator. The analogical character of this extension --- noted at \S 1.1 in the qualification on perceptual \textit{kinēsis} --- holds at every node of the chain at which the three-factor structure recurs.} The transition is perhaps the decisive hinge of the entire chain, since it is here that the physical world first acquires phenomenal determinacy for the ensouled body, and here that the material conditions for \textit{phantasia} (\gk{φαντασία} `imagination') and all subsequent cognitive operations are first laid down.

Three preconditions must obtain before $M_0 \rightarrow A_1$ can commence. First, the sensible object must be actively exercising its quality. Aristotle distinguishes carefully between what is ``capable of being seen'' and what is ``being seen'' (\textit{Metaphysics} IX.6, 1048b10--34): the former names the first actuality of the visible quality; the latter names the second actuality in which that quality is fully present to a perceiver. A colored surface in a lightless room possesses the first actuality of visibility without exercising it; only when illuminated and placed within an organism's perceptual field does the quality pass from dispositional to active presence. Second, the sense organ must hold its capacity in first actuality --- a structured readiness inherited from the organism's developmental history. §1.1 has already developed the three-grade schema of potentiality (bare capacity; developed capacity not currently exercised; exercise of capacity; cf. \textit{De Anima} II.5, 417a21--b2): the sense faculty at $A_1$'s dynamis-side stands at the second grade, the level of the \textit{hexis} (\gk{ἕξις} `settled disposition') that is poised for activation. The sense faculty in its first-actual state is not the bare potentiality of an infant who has never perceived but a structured readiness that has been informed by the organism's prior actualizations and is now disposed for the particular object that engages it. Third, the medium must be properly disposed to transmit the sensible form. Aristotle insists that direct contact between object and organ does not produce perception in the case of the senses that perceive at a distance.\footnote{For Aristotle's discussion of the various sensory media, see \textit{De Anima} II.7, 419a12--b2 (sight); II.8, 419b2--421a7 (sound and hearing); II.9, 421a8--422a7 (smell); II.10, 422a8--16 (taste); II.11, 422a8--424a15 (touch).} He develops the medium's role most explicitly at \textit{Sense and Sensibilia}: ``if the body which is heated or frozen is extensive, each part of it successively is affected by the part contiguous, while the part first changed in quality is so changed by the cause itself which originates the change'' (\textit{Sense and Sensibilia} 6, 447a3--7). The medium's affection is sequential. The sensible quality alters the part of the medium in contact with the object; that altered part alters the next part; and so on until the alteration reaches the sense-organ. Sound is genuinely successive, as Aristotle notes: we hear the clap of thunder after we see the flash because the alteration of the air takes measurable time to traverse the medium.\footnote{Light is the exception. Aristotle notes at \textit{Sense and Sensibilia} 6, 446b8--9 that light is transmitted instantaneously, actualizing the transparent in a single movement. The sequential character characteristic of sound and heat does not obtain for vision.}

The soul at $A_1$'s dynamis-side is not passive in the sense of being inert; it is receptive in the sense of being ontologically prepared. Heidegger's analysis of \textit{Geworfenheit} (`thrownness') provides a structural parallel: Dasein does not choose its being-attuned to the world but finds itself always already disposed in a particular way, affected by what it encounters before any reflective act of judgment (\textit{Being and Time} \S 29, H.135; Eng.~172--175). The structural point Heidegger isolates --- receptivity as structured readiness rather than as mere passivity --- holds equally for the Aristotelian sense-faculty at $A_1$'s dynamis-side: the organ is not a blank receptor awaiting impingement but a hylomorphically articulated capacity already poised to receive what falls within its proper range. The fuller existential structure of \textit{Befindlichkeit}\footnote{Heidegger's \textit{Befindlichkeit} is rendered throughout this dissertation as ``disposedness (\textit{Befindlichkeit})'', following Kisiel rather than Macquarrie \& Robinson's ``state-of-mind,'' with the italicized German parenthetical retained after every occurrence to preserve the technical anchor. \textit{Stimmung} is correspondingly rendered as ``attunement (\textit{Stimmung})'', following Rickert (\textit{Ambient Rhetoric}, 131, 146), rather than M\&R's ``mood.'' The reasoning is twofold. First, M\&R's ``state-of-mind'' carries mentalistic-internalist connotations that contradict Heidegger's explicit anti-internalist claim at \textit{Being and Time} \S 29, H.136: ``a mood [\textit{Stimmung}] assails us. It comes neither from the `outside' nor from the `inside', but arises out of Being-in-the-world, as a way of such Being.'' Kisiel's ``disposedness'' preserves the reflexive-receptive structure of \textit{sich befinden} (``finding-oneself-disposed-somehow'') without locating the disposing inside a subject. Second, Rickert's ``attunement'' for \textit{Stimmung} captures the ontological-atmospheric register Heidegger develops in \textit{The Fundamental Concepts of Metaphysics}: ``[\textit{Stimmung} is] like an atmosphere in which we first immerse ourselves in each case and which then attunes us through and through'' (GA 29/30, p.~67). The atmospheric-resonant character of attunement is particularly load-bearing for the dissertation's second main section, where I analyze human interaction with video-game and virtual-reality environments---contexts in which the ontological-immersive structure of being-attuned to a constituted environment is precisely what is at stake; M\&R's ``mood'' reads too psychologically for that analytical work. The glossary records the chosen renderings alongside M\&R's ``state-of-mind'' and ``mood'' for cross-reference. Quoted passages preserve the translators' English verbatim; readers may encounter ``state-of-mind'' or ``attunement'' for \textit{Befindlichkeit}, or ``mood'' for \textit{Stimmung}, in quoted material.}, which articulates the recursive disclosive character of attunement (\textit{Stimmung}), is reserved for the emotion section; for the present, the structural-readiness point suffices.

Gibson's ecological approach to visual perception reformulates the medium-precondition in a contemporary register. Where Aristotle speaks of the diaphanous as transparent and capable of receiving and transmitting the form of color, Gibson speaks of the ambient optic array as a structured field of light that specifies the surfaces and layouts of the environment (Gibson, \textit{The Ecological Approach to Visual Perception}, pp.~209--210). In both accounts, perception is not a matter of bare stimulus impinging on a passive receptor but a structured encounter between an already-organized environment and an already-capacitated organism. The tension is real and worth marking: Gibson is anti-mediation, holding that the structured optic array directly specifies environmental affordances without representational intermediary; Aristotle's hylomorphism, by contrast, is a doctrine of formal mediation, in which the medium transmits the form sequentially. The Gibsonian parallel is therefore restricted: it holds at the level of structural affordances --- the perceiver encounters an already-articulated environment, and the medium's structure is part of what makes that articulation available --- but it does not extend to full convergence with the Aristotelian doctrine of form-reception. With this restriction in place, the parallel is genuine and structurally useful, since it underwrites Aristotle's insistence that the perceptual situation is not the impingement of inert stimulus on an isolated receptor but a structured encounter constituted by object, medium, and organ in coordinated first actuality. $A_1$'s dynamis-side, accordingly, is not a single state but a constellation of preconditions, each in its proper first actuality before the perceptual motion $M_0 \rightarrow A_1$ can commence.\footnote{Daniel Gross's \textit{Uncomfortable Situations} extends the Gibsonian affordance framework into a humanities-side rhetorical-affordance theory that the present chapter's Aristotelian-hylomorphic reading anticipates: ``Rhetoric, as a humanistic form of affordance theory, posits a situation that is not neutral but is rather `persuasive,' threatening and promising with respect to our human being-in-the-world'' (Gross 22--23). The structured-encounter reading the present chapter develops of Aristotelian \textit{aisthēsis} --- in which the perceiver encounters an already-articulated environment and the medium's structure is part of what makes that articulation available --- is the cognitive-architectural condition for the kind of rhetorical-affordance dynamics Gross specifies at the contemporary register. Hyde's parallel reading develops the same structural point at the level of disclosure: \textit{aisthēsis} is a ``showing-forth'' (\textit{Aufzeigen}) of the world's significance to the perceiver (Hyde 89-92). Thomas Rickert's \textit{Ambient Rhetoric} (Pittsburgh, 2013) supplies a third Heidegger-grounded extension that explicitly ties Gibsonian affordance to dwelling-as-attunement: ``Rhetoric's action is material, affective, ecological, and emergent'' (Rickert 129), and ``Dwelling means the flourishing of life through attunement to the surrounding environment'' (Rickert 243). The resulting tripartite frame --- Gibsonian biological affordance, Gross's humanistic-rhetorical affordance, and Rickert's ambient-affordance --- articulates at the contemporary register the structured-encounter ontology that the present section Aristotelian-hylomorphic reading of \textit{aisthēsis} grounds at the cognitive-architectural level.}

\subsection{Hylomorphism and the Reception of Form}

The central and most controversial claim of Aristotle's perceptual theory is that the sense faculty receives the sensible form without the matter:\footnote{The scholarly literature on the form-without-matter doctrine is extensive. The principal interlocutors include Burnyeat, ``Is an Aristotelian Philosophy of Mind Still Credible? (A Draft),'' in Nussbaum and Rorty eds., \textit{Essays on Aristotle's De Anima} (Oxford: Clarendon, 1995), pp.~15--26; Sorabji, ``Body and Soul in Aristotle,'' \textit{Philosophy} 49 (1974): 63--89; and Caston, ``Aristotle and the Problem of Intentionality,'' \textit{Philosophy and Phenomenological Research} 57 (1997): 203--239. The section's task is to deploy the hylomorphic analysis as the structural basis for the dual-trace thesis developed in subsequent phases, not to adjudicate the controversy over whether the reception involves a literal physiological alteration (Sorabji) or a non-material spiritual reception (Burnyeat) or an intentional structure (Caston).}
\begin{adjustwidth}{0.5in}{0in}
Generally, about all perception, we can say that a sense is what has the power of receiving into itself the sensible forms of things without the matter, in the way in which a piece of wax takes on the impress of a signet-ring without the iron or gold; what produces the impression is a signet of bronze or gold, but not \textit{qua} bronze or gold: in a similar way the sense is affected by what is coloured or flavoured or sounding not insofar as each is what it is, but insofar as it is of such and such a sort and according to its form. (\textit{De Anima} II.12, 424a18--24)
\end{adjustwidth}
To receive the form without the matter is to undergo a change that is formally but not materially identical with the quality of the object; the eye, in seeing red, does not become materially red in the way that a dye stains cloth, but it does take on the formal structure of redness, becoming isomorphic with the quality it perceives. The signet-ring analogy is architecturally decisive. The ring impresses its engraved pattern into the wax; the wax takes the pattern; but the wax does not become iron or gold. What passes from ring to wax is the \textit{form} --- the structured design of the signet --- not the substrate that bears the form. Perception operates analogously: the color of the apple passes to the eye, the warmth of the flame to the skin, but the eye does not become apple-matter, nor the skin fire-matter. The sense and its organ are the same in fact but not in essence: the organ is a spatial magnitude, but the sensory power is a form or capacity in a magnitude rather than a magnitude itself (\textit{De Anima} II.12, 424a24--29).

The sensory power is a \textit{logos} (\gk{λόγος} `account, ratio') --- a proportion of the contraries proper to the sense's domain. Aristotle observes that sense is ``a sort of mean'' between the extremes of the sensible qualities (\textit{De Anima} II.11, 424a4--5) and develops the ratio analysis explicitly (III.2, 426a27--b8). Each sense is constituted by a proportion: the eye between light and dark, the ear between high and low pitch, the tongue between sweet and bitter. Perceptible qualities are themselves ratios of their underlying contraries, with grey as a specific proportion of white and black, or a flavour as a proportion of sweet and bitter. Perception is the organ's taking-on of the ratio that the perceptible quality instantiates: the sensory power adjusts its ratio to match the object's ratio, and this formal matching \textit{is} the perceptual act. The \textit{logos} framework explains the destruction condition that Aristotle posits: if the sensible exceeds the ratio the organ can receive, the motion acts on the organ as matter rather than as sense. Excess brightness destroys sight; excess loudness destroys hearing; excess heat destroys touch. When the stimulus intensity exceeds the proportional range of the organ, the organ is affected as matter rather than as sense, and the \textit{logos} itself is destroyed rather than realized (\textit{De Anima} III.2, 426b3--8). The motion is real, and it is genuine damage to the organ, but it is not perception. Perception requires that the stimulus fall within the ratio the organ can receive; excess falls outside the ratio, and the organ suffers rather than perceives. The destruction condition presupposes the hylomorphic analysis: if perception were merely material alteration, there would be no principled difference between perceiving and being damaged. The \textit{logos} is what makes that difference principled.

Heidegger's reading of Aristotelian perception in the \textit{Basic Concepts of Aristotelian Philosophy} radicalizes this analysis structurally. Heidegger reads Aristotle's account of \textit{aisthēsis} (\gk{αἴσθησις} `perception') as establishing that perception is itself a \textit{kritikon} (\gk{κριτικόν} `discerning, judging') --- ``a \gk{μέσον} [\textit{mesotēs} ``mean''] with the character of \gk{κριτικόν} [\textit{kritikon} ``discerning''], of the `ability-to-separate' one thing from another\dots It is a being-positioned toward possible objects, which is a \gk{δύναμις} [\textit{dynamis} ``potentiality''] in the sense of \gk{κριτική} [\textit{kritikē} ``faculty of judgment'']'' (\textit{BCAP}, p.~126). \textit{Aisthēsis} is not a passive imprinting but a discriminative event: the sense faculty, in receiving the form, simultaneously exercises its critical capacity to distinguish what it has received. Aristotle himself confirms the reading: the perceiving soul is said to ``make an affirmation or negation'' and to pursue or avoid accordingly (\textit{De Anima} III.7, 431a8--12). White's identification of \textit{to kritikon} as the unifying capacity that connects perception, \textit{phantasia}, and thought is the secondary literature's correlate of Heidegger's reading: White names \textit{to kritikon} ``the `distinguishing power,' collecting thought and sensation'' as against \textit{to kinetikon} (the locomotive power), with the major division at \textit{De Anima} III.9, 432a15--17 (White, ``Meaning of \textit{Phantasia},'' pp.~9--10), and reads \textit{nous} as ``a continuing and unifying theme throughout III, 3--11'' (White, p.~11, n.~39).\footnote{White's broader argument develops at pp.~9--11 the claim that \textit{to kritikon} is the architectural thread running through the discriminative powers of the soul --- perceiving, imagining, thinking --- such that the soul's capacity to discriminate underwrites both the formal reception of sensible qualities and the subsequent operations of \textit{phantasia} and \textit{nous}. The Heideggerian \textit{kritikon}-reading and White's \textit{to kritikon}-unification thesis converge at this structural site: discrimination is not a faculty added to perception but the formal character of perception's own activity, which then becomes the formal character of \textit{phantasia}'s and thought's activities as the chain proceeds.} The soul's capacity to discriminate --- that this is white and not black, sweet and not bitter --- presupposes that it has received the form of the sensible quality in a manner sufficiently determinate to ground such judgment. The hylomorphic reception of form is thus not a mere passive imprinting but a structurally discriminative event in which form-reception and discrimination are not two operations but one.\footnote{Hawhee extends the discriminative reading into a rhetorical-vision frame: ``Movement and life, \textit{kinēsis} and \textit{zōos}, are what confer rhetorical presence,'' and the eye's discriminative reception of the sensible quality is what makes possible the kind of ``rhetorical vision'' through which speech can move the audience by bringing the matter before their eyes (Hawhee 147--155). The structural relation is direct: the same \textit{kritikon}-structure that grounds Aristotelian \textit{aisthēsis} grounds the rhetorical capacity of the audience to receive \textit{phantasmata} from speech as discriminative-and-evaluative presentations.}

What persists after the perceptual act is over is not the sense faculty itself but the form the faculty has received, sustained in the organ's material substrate after the object has withdrawn. The persistence of this formal determination is the substrate on which \textit{phantasia} operates; \textit{phantasia} is the capacity that takes up the persisting trace and thereby makes thinking in the absence of objects possible. Persistence does not make \textit{phantasia} possible, since \textit{phantasia} is itself a developed capacity of the discriminative soul; \textit{phantasia} operates on persistence to make absence-presentation possible. The hylomorphic character of the original reception ensures that the persisting trace preserves the formal structure of the perceived quality and that the trace is therefore suitable as the material for \textit{phantasia}'s subsequent operations. Frede's analysis of the cognitive role of \textit{phantasia} contests a reductive reading of this persistence-character that she argues has been wrongly attributed to Aristotle: she notes that \textit{phantasiai} have ``sometimes been degraded to mere epiphenomena, the lingering after-images of sensations'' (Frede, \textit{The Cognitive Role of Phantasia in Aristotle}, p.~282) --- a characterization Frede herself rejects as failing to do justice to the full range of functions \textit{phantasia} performs in Aristotelian cognitive psychology (Frede, pp.~4--5). The resonant trace is not merely a fading echo but a structured residue retaining the formal character of the original perception while being materially reconstituted in the organ's ongoing physiological activity --- a point the next phases of the section develop.

\subsection{The Activity of Mover and Moved as One}

Aristotle's account of the perceptual act turns upon a principle inherited from the analysis of motion at $A_0$: the activity of the mover and the activity of the moved are one and the same activity, though they differ in their being. §1.1 established the principle at the motion-level: ``motion is in the movable. It is the fulfilment of this potentiality by the action of that which has the power of causing motion; and the actuality of that which has the power of causing motion is not other than the actuality of the movable; for it must be the fulfilment of \textit{both}'' (\textit{Physics} III.2, 202a13--20). The mover's \textit{poiēsis} (\gk{ποίησις} `making, producing') and the moved's \textit{pathēsis} (\gk{πάθησις} `being-affected, undergoing') are distinguished in account but not in substrate, ``one in substrate, different in being'' (cf. \textit{Physics} III.3, 202a21--24, on agent and patient; 202b13--14 on the Thebes-Athens analogy: one road, two descriptions).

Aristotle applies this principle directly to perception: ``the activity of the sensible object and that of the sense is one and the same activity, and yet the distinction between their being remains'' (\textit{De Anima} III.2, 425b26--27). The sensible object's exercise of its power (\textit{poiēsis}) and the sense faculty's reception of the form (\textit{pathēsis}) are not two events conjoined but one event under two \textit{logoi}. Aristotle illustrates with a striking example: ``a man may have hearing and yet not be hearing, and that which has a sound is not always sounding. But when that which can hear is actively hearing and that which can sound is sounding, then the actual hearing and the actual sound come about at the same time'' (\textit{De Anima} III.2, 425b28--426a1). The actual hearing and the actual sound are the same actuality under two descriptions; the sounding and the hearing arise together, in a single actualization.

The formula ``one in substrate, different in being'' is load-bearing across the section. It is introduced at the motion-level in \textit{Physics} III.2--3, redeployed at \textit{De Anima} III.2, 425b26--27 for perception, applied again at \textit{De Anima} III.7, 431a13--14 to the coordination of perception and appetite, and reapplied in this section's dual-trace thesis below. The formula's persistence across contexts reflects the chain's architectural unity: a single ontological pattern --- numerical identity under plural formal descriptions --- iterates across levels, each application specifying a new relation that the substrate-level unity secures. The audit-question the chapter must answer is whether each application is genuinely warranted by the cited Aristotelian text. At the motion-level and the perception-level, the textual warrant is explicit; at the perception-appetite level, Aristotle's own formulation at \textit{De Anima} III.7 (``the faculty of appetite and avoidance are not different, either from one another or from the faculty of sense-perception; but their being is different'') provides the warrant directly; at the dual-trace level developed in §F below, the formula is deployed as a structural extension consistent with the prior three applications.

The unity of mover and moved in the perceptual act establishes a fundamental principle of correlation: the sensible world and the sentient organism are not independent substances that happen to interact but correlative poles of a single actualization, each requiring the other for its full reality to be disclosed. Heidegger's existential reading sharpens this point in a different register. Dasein is not a static substance that subsequently enters into relations but a being constituted by its \textit{In-der-Welt-sein} (`being-in-the-world'); the receptivity that characterizes Dasein's being-in is not passivity but structured openness, in which what Dasein encounters is encountered as already articulated within an involvement-context (cf. \textit{Being and Time} \S 15--\S 18). The structural payoff for the present analysis is precise: receptivity, on the existential reading no less than on the Aristotelian, is not the inertness of an isolated receptor awaiting impingement but the structural openness of a comporting being whose mode of existence is to be already disposed toward what it can receive. The sense organ's first actuality (the dynamis-side of $A_1$) instantiates this structural openness at the level of perceptual capacity; Dasein's being-in-the-world generalizes the structure at the existential level. The recursive disclosive structure through which attunement (\textit{Stimmung}) conditions subsequent disclosure is reserved for the emotion section; the present analysis develops only the structural-openness point that the perceptual unity-thesis requires.

\subsection{Perception as Enmattered \textit{Alloiōsis} and the Four Senses of \textit{Pathos}}

Although the reception of form without matter is the defining characteristic of perception, Aristotle is careful to insist that the reception is not a purely formal event detached from the body. Perception is an enmattered \textit{alloiōsis} (\gk{ἀλλοίωσις} `alteration, qualitative change') --- a qualitative alteration that occurs in and through the material constitution of the sense organ. ``Sensation depends, as we have said, on a process of movement or affection from without, for it is held to be some sort of change in quality'' (\textit{De Anima} II.5, 416b33--417a1). The eye is physically affected by light; the ear by sound; the flesh is heated or cooled by the tangible object. These physiological changes are not mere accompaniments of perception but its material substrate; the formal reception could not occur without them. The soul, as Aristotle insists, is ``a substance in the sense of the form of a natural body having life potentially within it'' (\textit{De Anima} II.1, 412a19--22), and the activities of the soul --- including perception --- are inseparable from the body whose form it is.

Two kinds of \textit{alloiōsis} must be distinguished. The first is ordinary qualitative change, in which a substance's quality is altered by the imposition of a contrary: cold water becomes hot, a living body becomes diseased, a flame is extinguished. This kind of \textit{alloiōsis} is corrupting; the contrary supplants what was there. The second kind of \textit{alloiōsis} is preservative rather than corrupting. ``When we use the term `to be acted on,' we do not use it in one single sense, but in two: one is the destruction of one of two contraries by the other, the other is the preservation of what is potential by what is actual and like it'' (\textit{De Anima} II.5, 417a21--b2). The geometer learning geometry, the builder acquiring the art of building, the sense-organ receiving a sensible form: these are alterations that do not destroy the subject's prior state but actualize a potentiality already present as first actuality. Perception belongs to the preservative kind. When the eye sees red, the eye is not destroyed or replaced by red; its capacity to see is actualized, and the actualization is a fulfilment of that capacity rather than its negation. Aristotle is explicit: ``in the case of sense clearly the sensitive faculty already was potentially what the object makes it to be actually; the faculty is not affected or altered'' in the corruptive sense (\textit{De Anima} III.7, 431a3--5). The wax-signet analogy clarifies the point: the wax is altered by the ring --- it acquires a pattern it did not previously have --- but the alteration is not the wax's corruption; the wax is preserved as wax, and the pattern it receives is a determination of the wax rather than a substitution for it. The deeper structural register here is \textit{paschein} (\gk{πάσχειν} `being-affected'), under which Aristotle himself distinguishes two senses. On the one hand, \textit{paschein} expresses ``the extinction of one of two contraries by the other''; on the other, ``the maintenance of what is potential by the agency of what is actual and already like what is acted upon, as actual to potential'' (\textit{De Anima} II.5, 417b2--7). In one sense being-affected is a destruction; in the other it is a preservation in which a potential is brought into its actuality. Hearing a sound is a being-affected of the second sort: hearing is brought to its full being, not destroyed. The same structural point governs the perceptual \textit{alloiōsis} (cf. the parallel development at the level of higher-order \textit{pathē} in the emotion section, where the preservative-\textit{paschein} structure is deployed for the determinate emotions).

The structural register at issue is \textit{paschein} (\gk{πάσχειν} `being-affected'), and the perceptual \textit{alloiōsis} is the perceptual instance of \textit{paschein} in its preservative mode. The full structural frame within which this mode operates is given by Aristotle's fourfold articulation of \textit{pathos} (\gk{πάθος} `affection, being-affected') in \textit{Metaphysics} Δ.21:
\begin{adjustwidth}{0.5in}{0in}
We call an affection (1) a quality in respect of which a thing can be altered, e.g. white and black, sweet and bitter, heaviness and lightness, and all others of the kind. (2) The already actualized alterations. (3) Especially, injurious alterations and movements, and, above all, painful injuries. (4) Experiences pleasant or painful when on a large scale are called affections. (\textit{Metaphysics} Δ.21, 1022b15--21)
\end{adjustwidth}
The four senses articulate \textit{pathos} as the structure of being-occurrable-to that belongs to the way of being of a finite living being. Heidegger's reading of the fourfold in the \textit{Basic Concepts of Aristotelian Philosophy} treats the progression as a ``progressive narrowing,'' moving from the most general ontological structure (alterability) through the actual occurrence of alteration (\textit{energeia} of alteration) to the qualitative determination of that occurrence as harmful, to the magnitude of the harm. Heidegger gathers the four under a single thesis: ``From these four meanings, the \textit{genuine relatedness} of \gk{πάθος} becomes visible; it is related to the \textit{being of living things}, which is characterized by a thus-finding-oneself-again-and-again. The occurring to one befalls and strikes one in this disposition'' (\textit{BCAP}, p.~132). Heidegger then specifies that the occurrence need not without qualification have the character of the harmful: ``Aristotle recognizes a \gk{μεταβολή} [\textit{metabolē} ``change, alteration''], \gk{κίνησις} [\textit{kinēsis} ``movement''], \gk{ἀλλοίωσις} [\textit{alloiōsis} ``alteration, qualitative change''], in which \gk{πάσχειν} [\textit{paschein} ``being-affected''] has the character of \gk{σωτηρία} [\textit{sōtēria} ``preservation'']'' (\textit{BCAP}, p.~132).\footnote{Aristotle himself: ``Also the expression `to be acted upon' has more than one meaning; it may mean either the extinction of one of two contraries by the other, or the maintenance of what is potential by the agency of what is actual and already like what is acted upon, as actual to potential'' (\textit{On the Soul} II.5, 417b2--5).} There is not just one kind of \textit{paschein}: one is destruction by an opposite, the other is preservation by the actuality of what one was potentially. \textit{Pathos}, in its widest ontological reach, includes the structure by which a living being is brought into its proper actuality through its encounter with the world. Each of the four senses, on Aristotle's own analysis, is satisfied at the perceptual level: (1) alterability is the constitutive feature of the perceptual \textit{dynamis} (\textit{De Anima} II.5, 416b33--417a20); (2) actual alterations are realized in the resonant \textit{aisthēma} as the \textit{energeia} of the \textit{aisthētikon} at perceptual actualization; (3) harmful alterations are operative at the limit, where excess destroys the \textit{logos} (\textit{De Anima} III.2, 426a30--b8); (4) magnitudes are the structural feature by which perception admits of intensity and gradation. The four senses articulate the structural register within which the perceptual \textit{paschein} operates as a preservative mode of being-affected, with the destructive mode operative only at the limit. The fourfold and the preservative-\textit{paschein} distinction are taken up at the level of higher-order \textit{pathē} in the emotion section, where the structural prior commitments laid down here are deployed as the architecture of the determinate emotions; here the fourfold serves as the frame for the perceptual case alone.

Two consequences follow for the chain. First, the perceptual \textit{alloiōsis} is enmattered through and through: the formal reception is realized in and as the physiological alteration of the organ. Nussbaum notes the significance of this enmattered character when she observes that the activities of \textit{nous} in a living body ``must be linked with those of the faculty of \textit{phantasia}, which it would not be implausible to see as realized in some physiological occurrences'' (Nussbaum, \textit{The Role of Phantasia in Aristotle's Explanation of Action}, p.~238). If even \textit{noetic} activity requires physiological realization through \textit{phantasia}, then \textit{a fortiori} perception, which is the causal origin of \textit{phantasia}, must be thoroughly enmattered. The Heideggerian counterpart anticipates the emotion section's central thesis: the \textit{pathē} are the ground out of which speaking arises (BCAP, p.~176). Second, the preservative character of the perceptual \textit{paschein} ensures that the form received leaves a trace: the organ, having been formally altered, does not return to its pre-perceptual state the moment the object withdraws. Something persists. What persists, and how it does so, is the question the next phase addresses.

\subsection{Resonant \textit{Kinēsis}: The Persisting Impression}

The completion of the perceptual act does not terminate the \textit{kinēsis} initiated by the sensible object; it deposits a persisting impression in the ensouled body. Aristotle makes the point with striking consistency. In \textit{De Anima} he observes that ``even when the sensible objects are gone the sensings and imaginings continue to exist in the sense-organs'' (\textit{De Anima} III.2, 425b23--26). In \textit{On Dreams} he elaborates: ``The objects of sense-perception corresponding to each sensory organ produce sense-perception in us, and the affection due to their operation is present in the organs of sense not only when the perceptions are actualized, but even when they have departed'' (\textit{On Dreams} 2, 459a24--28). The residual affection is not a metaphor for an after-effect but the continuation of the perceptual motion in the organ after the originating cause has withdrawn. Aristotle illustrates the mechanism with a physical analogy: ``that part which has been heated by something hot, heats the part next to it, and this propagates the affection onwards to the starting-point. This must therefore happen in sense-perception, since actual perceiving is a qualitative change'' (\textit{On Dreams} 2, 459b1--6). The water, having been heated, does not instantaneously return to its pre-heated state when the fire is withdrawn; the warmth persists as a residual motion in the water's material substrate. So too in perception: the organ, having been formally altered, retains the alteration as a residual motion that persists for some time after the object has withdrawn. Aristotle gives a name to these residual affections: ``even when the external object of perception has departed, the impressions it has made persist, and are themselves objects of perception'' (\textit{On Dreams} 2, 460b2--3). The Greek is \textit{aisthēmata} (\gk{αἰσθήματα}) --- residual sense-impressions, the material of dreams, after-images, imagination, and the beginnings of memory.

This section adopts the term \textit{resonant kinēsis} for the residual motion understood as a whole --- the ongoing activity in the sensory apparatus that persists after the sensible object has withdrawn. The coinage marks two features of the trace that Aristotle's own lexicon does not consolidate into a single technical vocabulary. First, the trace is \textit{kinēsis}: a motion, not a static imprint, with the ongoing character of all Aristotelian \textit{kinēsis}. Second, it is resonant: it continues reverberating within the organ like sound continuing in a stilled bell, retaining the formal structure of the original motion while operating in the absence of its originating cause. White captures the reverberating character precisely: ``This lingering, resonating, echoing presence of sensible forms freed from their original matter appears to be the essence of \textit{phantasia} for Aristotle. The movement in the sense-power caused by the sensible object itself sets up a second movement which continues after the sensation has ceased and the object is gone'' (White, \textit{The Meaning of Phantasia in Aristotle's De Anima III, 3-8}, p.~498). The resonant \textit{kinēsis} is not the same in kind as the original perceptual act; it is a continuing exercise of the same formal structure, sustained by the organ's material substrate in the original cause's absence.\footnote{Caston's structural account of \textit{phantasia} as ``an echoing of the initial stimulation in the sense organs: a side-effect, like the original stimulation in character, but unable to compete with fresh, incoming stimulations'' (Caston 47, citing \textit{De Insomniis} 2, 459a23--b23) supplies the canonical contemporary anchor for the \textit{resonant kinēsis} coinage. Indeed, the \textit{resonant kinēsis} and \textit{phantasia}, for that matter, ``is a feeble sort of sensation'' (\textit{Rhetoric} I.11, 1370a28), and, as such, will be overridden by direct stimuli; this perceptual structuring, if you will, is what the present section generalizes to the dual-trace thesis: the residue carries both formal-imagistic content (\textit{resonant aisthēma}) and affective-valence content (\textit{resonant epithymia}), each subject to the same dynamic competition with fresh stimuli, each available for subsequent re-engagement by \textit{phantasia}.}

A terminological distinction must be made explicit before the dual-trace thesis is developed. The chapter operates with two related but non-synonymous terms: \textit{epithymia}\footnote{Aristotle defines \textit{epithymia} (\gk{ἐπιθυμία}) at \textit{De Anima} II.3, 414b5--6 as ``\textit{orexis} of the pleasant'' (\gk{τοῦ γὰρ ἡδέος ὄρεξις αὕτη})---the species of \textit{orexis} (genus) that all sense-perceiving animals possess insofar as they have touch + pleasure/pain (cf.\ 413b22--25; 414b1--6 distinguishing \textit{epithymia} / \textit{thymos} / \textit{boulēsis} as the three species of \textit{orexis}). I use \textit{epithymia} in this dissertation specifically for the basic hedonic valence (good/bad, pursue/avoid) co-given with live \textit{aisthēsis} at the present-object phase. Aristotle's broader applications of \textit{epithymia} to objects no longer present---e.g.\ \textit{Rhet.} I.11, 1370a25--b18 on the present-felt pleasure of remembering past drinks during fever-thirst---I treat as the re-engagement of \textit{resonant epithymia} (the past basic-valence-content preserved in the residual \textit{kinēsis}) producing live \textit{epithymia} in the present; the articulationally-concrete parallel---re-engagement of past \textit{pathē}-content preserved in memory producing a live \textit{pathē}---I name \textit{resonant pathē}.} and \textit{resonant epithymia}\footnote{I extend Aristotle's analyses of residual \textit{kinēsis}---which treat the formal-imagistic trace (\textit{phantasma} as \textit{eikōn / typos}, per \textit{On Memory} 450a26--b11; \textit{On Dreams} 459a24--460b3) as what persists in the sense-organs after the object departs---to claim that the residual \textit{kinēsis} also carries an affective-valence dimension I call \textit{resonant epithymia}. \textit{Mem.} 453a14--32 (the kinetic persistence of already-excited anger and terror) supplies partial textual support for affective-kinetic persistence; the dual-trace thesis generalizes this to the dormant-residual mode of basic hedonic tonality. The pair \textit{epithymia} (live) / \textit{resonant epithymia} (residual) follows Aristotle's own temporal contrast at \textit{On Dreams} 460b3--16 between live, \textit{enargēs} phantasia-perception and residual \textit{phantasmata} of dreams.}. \textit{Epithymia} names the perceptual event itself in its affective aspect --- the basic-valence dimension of the perceptual actualization as an event, the pleasant-or-painful that is co-given with every perceiving. \textit{Resonant epithymia} names the trace of that \textit{epithymia} persisting after the object's withdrawal --- the basic valence as the dispositional residue deposited by the perceptual event in the organ. The two are related as actuality to residue: the perceptual \textit{epithymia} is the event; the \textit{resonant epithymia} is the trace the event deposits. This distinction matters for what follows because the dual-trace thesis says that what persists in the organ is not a single undifferentiated impression but a composite trace with two formally distinct aspects, and the \textit{resonant epithymia} is one of those two aspects. The terminology is reserved this way: \textit{epithymia} for events, \textit{resonant epithymia} for the appetitive trace within the dual-trace composite. Together they are the two temporal phases of what \S 1.4 canonicalizes as \textit{basic affective valence}.

Jakob Von Uexküll's concept of \textit{Umwelt} offers a structurally productive analogue for the resonant \textit{kinēsis}'s ongoing function. The \textit{Umwelt} is not a passive backdrop against which perception occurs but a meaningful environment that the organism constitutes through its perceptual and effectual organs (von Uexküll, \textit{A Foray Into the Worlds of Animals and Humans}, pp.~222--224). The resonant \textit{kinēsis} can thus be understood as a modification of the organism's \textit{Umwelt}: the persisting impression does not merely linger in the organ as a static trace but actively shapes the organism's subsequent encounters with its environment, determining what it notices, what it ignores, and how it responds. Von Uexküll's framework is biological and anti-anthropocentric; the present analysis restricts the analogy to the basic-valence level, where it holds equally for rational and non-rational animals --- the persisting trace shapes the perceptual world the animal inhabits, whether or not that animal is capable of the higher-order operations \textit{phantasia} and \textit{doxa} will enable. The extension of the \textit{Umwelt}-analogy to rational-animal perception, where the soul's discriminative capacity and its access to articulated \textit{doxa} alter the structure of the perceptual world, requires further qualifications the present section does not undertake.\footnote{Gross's \textit{Uncomfortable Situations} develops a contemporary articulation of the \textit{Umwelt}/affordance framework in a directly relevant register: the human environment is constituted by ``negative'' and ``hostile'' affordances (Gross's signature extension of Gibson) that emerge precisely because the human \textit{Umwelt} is suffused with linguistic and rhetorical meaning. ``\textit{Rhetoric}, as a humanistic form of affordance theory, posits a situation that is not neutral but is rather `persuasive,' threatening and promising with respect to our human being-in-the-world'' (Gross 22-23). The extension to rational-animals can be located at the site where rhetorical-historical situatedness converts the bare biological \textit{Umwelt} into a rhetorical-affordance environment. Thomas Rickert's \textit{Ambient Rhetoric} supplies a Greek-classical anchor for the same structural insight at the register the present chapter primarily operates within: the pre-Socratic and Platonic \textit{periechon} --- ``that which surrounds, encompasses, or makes room for'' --- appears alongside \textit{chora} in Plato's \textit{Timaeus} and was already in play with Anaximenes, Hippocrates, and Empedocles, indexing an awareness that the encompassing environment is itself alive and that it constitutes the human condition both materially and spiritually (Rickert 6--7). Rickert's recovery culminates in the claim that taking rhetoric as ambient amounts to ``a renewed sensibility to what the ancient Greeks called \textit{periechon}'' (Rickert 285). The \textit{Umwelt}-as-modern-biological and \textit{periechon}-as-ancient-Greek framings together close the temporal arc the present analysis traces from the Aristotelian medium-doctrine through contemporary affordance theory, anchoring the chapter's structurally productive analogue at the Greek-classical register the dissertation primarily inhabits.}

The signet-ring analogy, extended to three phases, illuminates what the resonant \textit{kinēsis} adds to the hylomorphic account. (1) \textit{Potentiality}: The wax is pliable and smooth; the ring is held above. The sense-organ is ready and the sensible object nearby but unseen --- $A_1$'s dynamis-side is complete, $A_1$'s actualization not yet underway. (2) \textit{Actualization}: The ring presses into the wax; the ring and wax are jointly actual in the impressing; the actual hearing and the actual sound come about at the same time. This is $A_1$ itself: the perceptual act as the shared actuality of sensible object and sense faculty. (3) \textit{Residual trace}: The ring is withdrawn, the actualization proper ends --- but the wax has undergone preservative \textit{alloiōsis}, and the impression persists. The \textit{aisthēmata} remain. The resonant \textit{kinēsis} is the organ's continuing exercise of the form received, sustained in the object's absence by the organ's own material substrate. An alteration that vanished with its cause could not ground further actualities; a completed actuality that terminated in itself could not serve as the unmoved originator of the next motion. But the resonant \textit{kinēsis} persists, and in persisting it becomes the material upon which the \textit{phantastic} motion $M_1 \rightarrow A_2$ will operate. The character of that residue --- whether it is a single undifferentiated impression or a structured composite --- determines what the subsequent stages can generate from it. The next phase argues that the residue is not single but dual.

\subsection{The Dual-Trace Thesis: Resonant \textit{Aisthēma} and Resonant \textit{Orexis}}

The persisting impression deposited by the completed perceptual act is not a single, undifferentiated residue but a dual trace comprising two distinct yet inseparable moments: the resonant \textit{aisthēma} (the persisting sensory form) and the \textit{resonant epithymia} (the persisting appetitive valence). Every perception is simultaneously a discrimination of a sensible quality and an evaluation of that quality as pleasant or painful, beneficial or harmful; accordingly, every perception deposits in the ensouled body both a formal trace of what was perceived and an affective trace of how what was perceived was evaluated.\footnote{Nussbaum's reading of \textit{phantasia}'s structural role specifies the cognitive-architectural condition for the dual-trace thesis: \textit{phantasia} ``ties abstract thought to concrete perceptible objects or situations'' (Nussbaum 265). Caston's contemporary articulation supplies the structural mechanism: ``By making \textit{phantasia} (rather than belief) the common currency among mental states, Aristotle emphasizes a form of intentionality which is more basic than the conceptual, and firmly rooted in the general character of perceptual experience'' (Caston 52). The dual-trace's formal-and-affective co-given character is, on these readings, not an \textit{ad hoc} addition but the structural consequence of \textit{phantasia}'s placement as the common pre-conceptual currency through which both cognitive and conative content are carried.} The dual trace is one in substrate --- a single enmattered affection in the organ --- and two in being, corresponding to the perceptual and appetitive faculties' distinct ways of being.

The argument proceeds along two converging lines. The first is analytic. Aristotle establishes a chain of entailments in \textit{De Anima}:
\begin{adjustwidth}{0.5in}{0in}
If any order of living things has the sensory, it must also have the appetitive; for appetite is the genus of which desire, passion, and wish are the species; now all animals have one sense at least, viz. touch, and whatever has a sense has the capacity for pleasure and pain and therefore has pleasant and painful objects present to it, and wherever these are present, there is desire, for desire is appetition of what is pleasant. (\textit{De Anima} II.3, 414b1--6)
\end{adjustwidth}
Every act of sensation, in this chain, analytically entails affective valence. To sense at all is already to be oriented toward the pleasant and away from the painful; to have pleasant and painful objects present is to have the appetite those objects engage. Sensation is not a neutral registration of qualities that can then be evaluated by some further faculty; sensation is already evaluative in its very structure. Aristotle develops the coordination of perception and appetite directly:
\begin{adjustwidth}{0.5in}{0in}
To perceive then is like bare asserting or thinking; but when the object is pleasant or painful, the soul makes a sort of affirmation or negation, and pursues or avoids the object. To feel pleasure or pain is to act with the sensitive mean towards what is good or bad as such. Both avoidance and appetite when actual are identical with this: the faculty of appetite and avoidance are not different, either from one another or from the faculty of sense-perception; but their being is different. (\textit{De Anima} III.7, 431a8--14)
\end{adjustwidth}
Perception and appetite are numerically one in substrate --- one faculty, one event --- but they differ in being. When the object is pleasant or painful, the perceptual act is simultaneously an affirmation or negation, a pursuit or avoidance. There is not first a perception and then, by a separate cognitive operation, an evaluation; the perception itself carries the evaluation as a formal aspect of its being. The residual motion that persists after the perceptual act must inherit this dual structure. The sensation that deposited the trace was itself numerically one and two in being; its \textit{pathos} was formally both perceptual and appetitive. Unless one can argue that the residue discards one aspect while retaining the other --- which Aristotle's hylomorphism does not permit --- the full \textit{pathos} persists, with both formal aspects intact. The enmattered-accounts argument from \textit{De Anima} I.1 supports the inheritance directly. The affections of soul are ``enmattered accounts'' (\textit{logoi enyloi}) whose definition must combine formal and material aspects; the physicist gives the material account (a boiling of the blood around the heart), the dialectician the formal account (a desire for retaliation), and the full definition requires both (\textit{De Anima} I.1, 403a25--b19). If the residual motion is an enmattered \textit{pathos}, and the \textit{pathos} has two formally distinct aspects (perceptual and orectic), then the residue preserves both aspects as formal aspects of a single material substrate.

The second line of argument is a \textit{modus tollens}. If the residual motion were exhausted by its formal-epistemic aspect --- if what persisted were only the resonant \textit{aisthēma} --- then mere retention of sensible form would suffice to originate pursuit and avoidance. The soul's knowledge \textit{that the object is a bear} would, by itself, move the animal to flee. But Aristotle insists that bare perception is ``like bare asserting or thinking''; it is only \textit{when the object is pleasant or painful} that the soul ``makes a sort of affirmation or negation, and pursues or avoids the object'' (\textit{De Anima} III.7, 431a8--14). The formal content alone is motivationally inert. A distinct operation --- the hedonic-orectic valuation of the object as pleasant or painful --- is required to initiate pursuit or avoidance; the source of all movement, as Aristotle states, is appetitive: ``that which originates movement must be specifically one, viz. the faculty of appetite as such (or rather farthest back of all the object of that faculty; for it is it that itself remaining unmoved originates the movement by being apprehended in thought or imagination)'' (\textit{De Anima} III.10, 433b10--14).\footnote{Caston's reading of the pleasure/pain and appetite/desire connection to perception further strengthens my Aristotelian reading that every moment of sensation is co-given with a hedonic tonality: ``any thing which has at least one sense can feel pleasure and pain, and if so, will have appetite and hence desire'' (Caston 22, citing \textit{DA} 2.2, 413b22-23; 2.3, 414b4-5; 3.11, 434a2-3). The co-givenness of pleasure and pain with sensation is therefore not an add-on Aristotelian doctrine but the structural-functional condition for the perceiving animal's having appetite at all; the dual-trace's appetitive aspect is the residue of this constitutive co-givenness. See: \textit{De Anima} III.10, 433b16-19: ``that which moves without itself being moved is the realizable good, that which at once moves and is moved is the faculty of appetite (for that which is moved is moved insofar as it desires, and appetite in the sense of actual appetite \textit{is} a kind of movement).''} The motivational autonomy of the appetitive trace is therefore not an ad hoc addition to the hylomorphic analysis but a structural requirement of Aristotle's own commitment to appetite as the source of animal movement.

What the \textit{resonant epithymia} preserves is what the chapter names the \textit{hedonic tonality} of the perceptual event --- the pleasant-or-painful aspect that is co-given with every perceiving and that persists as a dispositional trace after the object withdraws. The term is grounded in Heidegger's reading of Aristotelian \textit{hēdonē} (\gk{ἡδονή} `pleasure') and \textit{lypē} (\gk{λύπη} `pain, distress') in the \textit{Basic Concepts of Aristotelian Philosophy}. Heidegger insists, first, that \textit{hēdonē} is what it is ``in the moment,'' not in time as a determinate span: \textit{hēdonē} ``is what it is `in the moment,' \gk{μὴ ἐν χρόνῳ}, `not in time' in the sense of a determinate span,'' and while \textit{hēdonē} is itself a kind of \textit{kinēsis}, it is not a movement in the sense of a process spanning time but ``a \textit{determination of the presentness of being-there as such}'' (BCAP, p.~164). Second, Heidegger insists that \textit{hēdonē} and \textit{lypē} are co-given with every mode of being-there: ``with every \gk{πάθος} (\textit{pathos}), but equally with every perceiving, every thinking, considering, with \gk{θεωρία} (\textit{theōria}), to the extent that they are basic modes of living, \gk{ἡδονή} (\textit{hēdonē}) is an inseparable companion'' (BCAP, p.~166). There is no mode of being-there from which pleasure-or-pain is absent: every act of perceiving, judging, contemplating, and suffering is tonally colored by \textit{hēdonē} and \textit{lypē}. Hedonic tonality, on this reading, is not an add-on to perception that arrives after the perceptual event is complete but a constitutive aspect of the perceptual event itself. The \textit{resonant epithymia} is the persisting trace of this constitutive aspect. What persists after perception is therefore not merely the formal structure of the perceived quality but the formal structure together with the hedonic tonality under which the quality was apprehended.

The dual-trace thesis delivers the structural payoff for the same-chain-identity claim established in §1.1. The chain runs through perception $\rightarrow$ \textit{phantasma} $\rightarrow$ desire-actualization not because two parallel chains have been collapsed into one but because perception itself, in its very structure, deposits both a formal trace and an appetitive trace. The basic appetitive valence is co-given with perception itself; there is no perceptual stage without a desire-component. The perception-to-movement chain IS the actualization-of-desire chain not by stipulation but because the perceptual \textit{epithymia} analytically entails the appetitive aspect that drives subsequent action. Same-chain identity is therefore not a high-level claim imposed on the chain from outside but the structural articulation of what the dual-trace thesis establishes at its core: the formal and appetitive moments of perception are one in substrate, and the chain that runs through both is one chain.

A distinction at the close of the phase orients the analysis forward. The dual trace deposited at $A_1$ is basic valence in its perceptual instance --- the \textit{resonant epithymia} that preserves the pleasant-or-painful aspect of the perceptual event. The higher-order \textit{pathē} that Aristotle catalogs in the \textit{Rhetoric} --- anger as ``desire accompanied by pain for a conspicuous revenge for a conspicuous slight'' (\textit{Rhetoric} II.2, 1378a30--32), fear as ``pain or disturbance due to imagining some destructive or painful evil in the future'' (II.5, 1382a20--22), pity as a ``feeling of pain at an apparent evil\ldots\ which we might expect to befall ourselves or some friend of ours'' (II.8, 1385b15--17) --- are not at issue at this stage. The higher-order \textit{pathē} arise only downstream of \textit{doxa}'s orthogonal ratification of $A_3$ content, where the cognitive apparatus articulates the perceptual basic valence into the determinate evaluative concretions the emotion section develops. The structural distinction between \textit{epithymia} as basic valence (operative at $A_1$ onward) and the determinate higher-order \textit{pathē} (arising only downstream of \textit{doxa}'s ratification) is load-bearing for the chapter and is the central thesis of the emotion section; here it appears as the closing structural finding of the perceptual analysis.

\subsection{$A_1$ as Ground of $A_2$}

The completed perceptual act, together with the dual residual trace it deposits, constitutes $A_1$ --- the completed perception of the chain and the ground upon which the next stage of actualization, $A_2$, will be built. $A_1$ is not merely an endpoint but a threshold: it names the state of the ensouled body after perception has occurred, a state in which the organism carries within itself the formal and appetitive determinations of its encounter with the sensible world. The architectural role follows from the iterated chain structure established at $A_0$. Each completed actuality is simultaneously the terminus of its antecedent motion and the unmoved originator of its successor. $A_1$ is the terminus of the perceptual motion $M_0 \rightarrow A_1$; it is also the unmoved originator of the \textit{phantastic} motion $M_1 \rightarrow A_2$ through which the \textit{phantasma} proper will emerge. The resonant \textit{kinēsis} deposited at $A_1$ is the material on which \textit{phantasia} will operate at the next stage. The \textit{phantasma} is not created \textit{ex nihilo} from the sensible object's form; it is generated from the dual trace $A_1$ has left in the sensory apparatus. \textit{Phantasia}, as Aristotle defines it at \textit{De Anima} III.3, 428b10--17, is a \textit{kinēsis} arising from actual sensation and necessarily similar in character to the sensation from which it arose:
\begin{adjustwidth}{0.5in}{0in}
But since when one thing has been set in motion another thing may be moved by it, and imagination is held to be a movement and to be impossible without sensation, i.e. to occur in beings that are percipient and to have for its content what can be perceived, and since movement may be produced by actual sensation and that movement is necessarily similar in character to the sensation itself, this movement cannot exist apart from sensation or in creatures that do not perceive, and its possessor does and undergoes many things in virtue of it, and it is true and false. (\textit{De Anima} III.3, 428b10--17)
\end{adjustwidth}
The similarity is grounded in the resonant \textit{kinēsis}: the \textit{phantasma} inherits both the formal aspect (resonant \textit{aisthēma}, making the image recognizable as of that object) and the affective aspect (\textit{resonant epithymia}, making the image matter to the perceiver) from the residual motion deposited at $A_1$. Without the resonant \textit{aisthēma}, there would be no formal structure for \textit{phantasia} to preserve, recombine, or transform; without the \textit{resonant epithymia}, the \textit{phantasma} would be a mere formal schema --- a picture without significance --- rather than the charged image that drives desire, aversion, and deliberation.

The phantasia-foreshadowing landing now becomes explicit. §1.1 foreshadowed \textit{phantasia} as the temporal-narrative-synthesizing capacity through which a finite, embodied, comporting being cognizes time as a coherent narrative of past, present, and future --- the capacity through which memory, present perception, and anticipation are unified into a single phenomenal manifold. §1.2 deposits the material on which that synthesizing capacity operates. The resonant \textit{kinēsis} is not yet itself a temporal-narrative manifold; it is the formal-and-affective residue from which such a manifold will be constructed by \textit{phantasia}'s synthetic activity in subsequent sections. The chain's twofold articulation --- as the actualization of desire and as the perception-to-movement chain --- runs through this junction: the dual trace at $A_1$ is what makes the chain genuinely one (per the dual-trace defense of same-chain identity established in §F above) and what supplies \textit{phantasia} with the material for its temporal-narrative synthesis.

The threshold of $A_2$ can now be stated. The dual trace is in place. The resonant \textit{aisthēma} preserves the formal structure of the perceived quality; the \textit{resonant epithymia} preserves the hedonic tonality under which the quality was apprehended. The two are one in substrate and two in being, deposited as a single enmattered affection that the organ now carries. The next section takes up the \textit{phantastic} motion $M_1 \rightarrow A_2$ through which the dual trace is taken up by \textit{phantasia} and rendered as the \textit{phantasma} proper.

\end{document}